\part{图形与比较方式}


\section{正负双向条与误差标记}
\key{min} 默认是零。设置负的下限后，零线位于
$W(-\mathrm{min})/(\mathrm{max}-\mathrm{min})$，
负值向左、正值向右。非对称范围的零线不在轨道中心。
显示零点两侧时应让量程包含零；v0.0.8 的一般量程只要求下限小于上限。

\begin{example}{示例 10：共用零线的双向条}
\gradbarssetup{
  min=-50,max=50,
  width=30mm,
  negative color=lightred
}
\gradbar{-30}

\gradbar{0}

\gradbar{30}
\end{example}

未显式选择主题时，负值保留 \key{negative color=lightred} 的默认纯色。
显式设置 \key{theme=blue} 等主题后，负值采用该主题的互补色，填充方式与主题一致。
也可以为负值独立指定渐变：
\begin{example}{示例 10a：负值独立渐变}
\gradbar[min=-100,max=100,width=30mm,label width=3em,
  left color=teal,right color=teal!25,
  negative fill=gradient,
  negative left color=magenta!25,
  negative right color=magenta]{-60}
\end{example}
负值的 \key{left} 端是远离零点的一端，\key{right} 端靠近零点；\key{negative color}
仍是未指定渐变时的后备纯色，条件配色和 \key{print} 模式继续优先。

跨零范围不接受 \key{value format=percent}，
避免把位置比例当作有意义的百分比；百分数变化可直接使用 \key{unit}。

\begin{example}{示例 11：对称与非对称误差}
\gradbar[error=5]{72}

\gradbar[error minus=4,
  error plus=9]{72}

\gradbar[min=-50,max=50,
  error minus=8,error plus=12]{-5}
\end{example}

\key{error=e} 表示区间 $[v-e,v+e]$；
\key{error minus} 和 \key{error plus} 分别指定非负的上下误差幅度，
不是区间端点。只指定一侧时，另一侧默认为零。
误差单位与数据相同，\key{error color} 控制线色；
\key{error={}} 清除两侧误差。误差含义由用户说明，宏包不计算置信区间。

\begin{example}{示例 12：组合效果}
\gradbar[
  min=-50,max=50,width=30mm,
  rounded=2pt,target=20,
  error=5,label=end
]{-25}
\end{example}

\note{数值及误差端点超出范围时，默认发出警告并截断显示，
原始标签保持不变。截断的误差端点不代表真实端点。
严谨的区间图建议使用 \key{overflow=error}，或扩大范围以显示完整误差。}


\subsection{同轨呈现两个正负观测值}
\verb|\gradsignedbar[选项]{负值}{正值}| 由 core 内置，
\key{libraries=none} 即可使用。两个参数是独立观测值，不计算净值或差值。
负侧必须小于等于零，正侧必须大于等于零；量程须满足
$\mathrm{min}<0<\mathrm{max}$，也可使用非对称或命名共享量程。
\begin{example}{一条轨道、两侧数据及误差}
\gradsignedbar[min=-50,max=50,width=30mm,
  error=8,theme=blue]{-24}{24}

\gradsignedbar[min=-20,max=80,width=30mm,
  error minus=2,error plus=5,theme=teal]{-10}{40}

\gradsignedbar[min=-50,max=50,width=30mm,
  theme=accessible]{NA}{24}
\end{example}
公共轨道、零线与参考标记只绘制一次。误差参数分别作用于两个观测值，
误差区间可以跨过零点；越界仍按 \key{overflow} 处理，不因同轨而裁剪至某一侧。
每侧可独立使用 \key{NA} 或空参数表示缺失，零值表示零长度，不视为缺失。
默认在轨道外显示“负值 / 正值”，分别使用现有数字格式与单位；
支持 \key{label=outside}、\key{left}、\key{none}，\key{text} 可覆盖整个标签。
此命令只绘制普通条形，不接受其他复合图形模式；不增加原生表格列类型。
自动替代文本同时描述两个观测值；已有 TikZ 图内由外层图形提供描述。
自定义标量后端须显式设置 \key{backend fallback=tikz}。
